% ===== §9 The Three Beauties: Perception =====
\section{The Three Beauties and the Beauty of Perception}
\label{sec:three-beauties}

\noindent
The reduction chain has reached its terminus, the reframing of field-building as aesthetic education, and the coupling and its orthogonality guard have been set out. What remains, before the four-step cultivation can be developed in detail, is to lay out the structure of the capacity that the cultivation builds, and that structure has three parts. The capacity to build a generative field is at once a capacity to perceive, a capacity to practise, and a capacity to maintain, and these are not three separate skills but three faces of a single trained sensibility, distinguished by what they do in the cycle the field-building enacts. This section introduces the three and develops the first; the two that follow develop the others. The three together are the static anatomy of the capacity, the cross-section of what a field-builder has; the four-step model that comes later is the same capacity in motion, the cycle that the anatomy turns. It is worth holding both pictures, the anatomy and the cycle, because each shows what the other cannot, the parts and their turning.

\subsection{Three faces of one sensibility}
\label{sub:three-faces}

The capacity to build a generative field decomposes into three, by a division that follows the structure of any generative act. To build a field one must first perceive its state, register where it stands and what here would be appropriate; this is the beauty of perception, the input face, the discernment of the field as it is. One must then act upon it, perform the gesture, make the placement, that sends the path the generative way; this is the beauty of practice, the output face, the generation of appropriateness in the concrete act. And one must sustain this across time, keep the field accumulating through repetition rather than letting it dull or spin, return to it and turn it again without flattening; this is the beauty of maintenance, the temporal face, the holding of the good cycle in time. Perception discerns the direction, practice sends the path that way, maintenance keeps it turning, and the three are one sensibility because the same resolution upon the field's generative direction underlies all three, perceiving it, acting to extend it, and sustaining the action across cycles.

These map onto the inference and action of a single generative process, which is why they are faces and not separate faculties. Perception is the inferential moment, the registration of the field's state; practice is the active moment, the intervention that alters the state toward accumulation; and maintenance is the persistence of the inference-and-action across time, the keeping of the process turning rather than letting it run down. The unity matters because it forbids a mistake the tradition has made, the treating of aesthetic education as the cultivation of perception alone, the refinement of a taste that discerns without building or sustaining. A sensibility that perceived the field finely but could not act to build it, or could build it once but not sustain it, would not be the capacity this paper is about, which is the capacity to build generative fields and not merely to appreciate them. The three faces must be cultivated together, and their cultivation together is what the four-step model, taken up later, organises into a single turning cycle.

The order in which the three are presented, perception then practice then maintenance, is not arbitrary, and saying why it is the right order shows how the three depend on one another. Perception comes first because the other two presuppose it: one cannot act to build a field whose state one has not perceived, since the apt act is apt to a state, and one cannot sustain across time a field one cannot perceive to be flattening or gaining, since maintenance is the correction of a drift that must first be seen. Practice comes second because it presupposes perception and is presupposed by maintenance: it acts upon what perception has registered, and maintenance is the repetition across time of the perceiving-and-acting that practice, in a single instance, performs. Maintenance comes third because it is the temporal extension of the first two, the keeping-turning of a cycle whose single turn is a perceiving and an acting. The order is thus an order of dependence, each face presupposing the one before, and it is also the order of the cycle, perception opening onto practice and practice settling into the maintenance that returns to perception. But the dependence does not make the later faces less important than the earlier, and here the tradition's error must be guarded against from the other side. That perception comes first in the order of dependence does not make it first in importance, as though the others were mere applications of a primary discernment; on the contrary, as the section on practice will argue, in the intimate field it is practice that is dominant, the building on which the field's accumulation actually turns, and perception is subordinate to the practice it serves. The order of presentation follows dependence; the order of importance follows the field's state and is read dialectically; and the two orders must not be confused, on pain of reinstating, as a doctrine of perception's primacy, the very over-valuation of discernment that the unity of the three faces was meant to correct.

\subsection{The beauty of perception}
\label{sub:beauty-perception}

The beauty of perception is the discernment of the field's state, the registration of where it stands, whether it is accumulating or spinning, and what here, in this state, would be appropriate. It is the face of the sensibility that the coupling section analysed as resolution: a fine perception discriminates the field's generative direction at fine grain, catches the slight flattening before it compounds, registers the small contingency that a coarse perception would absorb without trace. To perceive the field beautifully is to see it as it is, undistorted by the prior that would pre-sort it and undimmed by the habit that would render it invisible, and the cultivation of this face is the raising of the resolution upon the field's state to the point where its accumulating or dissipating character is perceived immediately, as a trained ear perceives a discord. The detailed account of how this perception is possible, and of the traditions that contend over what seeing the field is, is the work of the section that follows this part; here the task is to place perception within the anatomy of the capacity and to draw the analysis of its principal and secondary aspect that the method requires.

\begin{casebox}[The coarse and the fine perception of one evening]
Two partners come home to the same evening, and the difference between a coarse and a fine perception of it can be made plain. The coarse perception registers only the gross states: the other is plainly angry, or plainly happy, or plainly fine, and below the threshold of the plain nothing is seen. To such a perception the evening presents itself as one of a few large categories, and the field is read only when its state has become unmistakable, which is usually too late to do anything but react to what has already compounded. The fine perception registers what the coarse one misses: the small tightness beneath the reported fineness, the slight flatness in a greeting that says the day went harder than the words admit, the faint lift that means something went unexpectedly well, the early and subtle signature of a field beginning to flatten while the flattening is still small enough to meet. The two perceptions face the same evening; one sees three or four states and the other sees a continuous and finely graded field, and the difference is the resolution the cultivation raises. It is the difference between a perception that reads the field only when its state has become gross and a perception that reads the field while its state is still subtle, which is the difference between meeting a drift while it is small and meeting it only after it has compounded into something no longer easily met.
\end{casebox}

\subsection{The principal and secondary aspect within perception}
\label{sub:perception-aspect}

The method announced that each contradiction has a dominant and a subordinate aspect, and that the two can transform into one another, and the analysis applies within the beauty of perception itself, where it concerns the contradiction between perceiving the field's whole state and perceiving the particular contingency within it. These are two moments of perception, the registration of the field as a standing condition and the catching of the singular unplanned thing that occurs within it, and in the perception of an intimate field they are not equally weighted. The catching of the contingency is the dominant aspect, because the field's state is most truthfully revealed in how it meets the unplanned, the contingency being the place where a flattening prior is most exposed and a living perception most tested; a perception that registered the standing state but missed the contingencies would be a perception of the field's surface and not its life. Yet the two transform into one another. When a relation has become so volatile that every moment is contingency and no standing state can be read, the dominant aspect shifts to the registration of the whole, the recovery of a sense of where the field stands beneath the churn, without which the contingencies cannot even be placed. The dominant aspect is the catching of contingency in the ordinary case and the registration of the whole in the volatile one, and which dominates is read from the field, as the method requires, rather than fixed in advance.

This is the first instance, set out in full so that the later instances can be brief, of the aspect-analysis the method promised, and it shows the analysis doing real work rather than decorating the exposition. It tells the field-builder where, within the work of perception, attention now belongs: ordinarily on the contingency, the unplanned thing that reveals the field's life, but in the volatile field on the recovery of the whole, without which the unplanned cannot be read. The same will hold, with different content, for practice and for maintenance, each carrying its own contradiction between a dominant and a subordinate aspect, each transforming as the field's state changes, and each to be read from the field rather than legislated. The anatomy of the capacity is thus not static after all, even in cross-section, because each of its faces carries within it a contradiction whose dominant aspect migrates, and the reading of those migrations is the standing exercise of the very sensibility the anatomy describes.
